\documentclass[../../main.tex]{subfiles}% 注意这里的文件路径不能用 ./main.tex ,否则用latexmk编译子文件会报错
\graphicspath{{\subfix{./image/}}{\subfix{../../image/}}} % 指定图片引用路径,后续可以直接使用图片文件名
% 如果图片存放在上级目录,则使用 ../../image/ 作为备用路径

% 例如:
% \begin{figure}[H]
% \centering
% \includegraphics[width=0.3\textwidth]{图.png}
% \caption{}
% \label{figure:图}
% \end{figure}
% 注意:上述\label{}一定要放在\caption{}之后,否则引用图片序号会只会显示??.

\begin{document}

\section{问题的提法}\label{section:ch07-问题提法}

本章考察 Cauchy 问题
\begin{align*}
\begin{cases}
\partial_t u+\partial_x[\phi(u)]=0 & \text{在 }\mathbb{R}\times(0,T)=:S,\\
u=u_0 & \text{在 }\mathbb{R}\times\{0\},
\end{cases}
\qquad(P)
\end{align*}
其中 $T>0$，$u_0$ 为 $\mathbb{R}$ 上的正有限 Radon 测度。关于 $u_0$ 为 $\mathbb{R}$ 上有限带号 Radon 测度的情形，参见文献~\cite{ref19}。

\subsection{假设与预备注记}\label{subsection:ch07-假设}

假设
\begin{align*}
u_0&\in\mathscr{R}_f^+(\mathbb{R}), \qquad\text{(A0)}\\
\phi&\in\mathrm{Lip}([0,\infty)),\qquad \lim_{y\to\infty}\frac{\phi(y)}{y}=0 \qquad\text{(A1)}
\end{align*}
（不失一般性可设 $\phi(0)=0$）。需指出，$\lim_{y\to\infty}\phi(y)/y=C_\phi\neq0$ 的一般情形可通过把 $x$ 换为 $x-C_\phi t$ 处理（细节见文献~\cite{ref18}）。由 (A1)（及条件 $\phi(0)=0$）存在 $M>0$，使
\begin{align}
|\phi'(y)|\leqslant M\ \text{对几乎所有 }y>0,\qquad |\phi(y)|\leqslant My\ \text{对每个 }y>0. \label{align:ch07-7.1}
\end{align}

将 $u_0$ 按 Lebesgue 测度分解为 $u_0=u_{0r}\,\mathrm{d}x+u_{0s}$，其中 $u_{0r}\in L^1(\mathbb{R})$ 为绝对连续部分的密度，$u_{0s}\in\mathscr{R}_f^+(\mathbb{R})$ 为奇异部分（见\cref{definition:ch01-Lebesgue分解}）。我们还将使用下列更强的假设（见\cref{section:ch07-唯一性}）：
\begin{align*}
&\text{奇异部分 }u_{0s}\text{ 为 Dirac 质量的有限叠加，}\\
&u_{0s}=\sum_{j=1}^p c_j\delta_{x_j}\qquad(x_1<x_2<\cdots<x_p;\ c_j>0,\ 1\leqslant j\leqslant p). \qquad\text{(A0')}
\end{align*}

为特定目的，我们再作如下假设（见\cref{section:ch07-存在性}）：
\begin{align*}
\begin{cases}
\phi\in C^1([0,\infty)),\ \text{且对每个 }y_1>0,\text{ 存在 }a,b\geqslant0,\ a+b>0,\\
\text{使 }\phi'\text{ 在 }[y_1-a,y_1+b]\cap[0,\infty)\text{ 上严格单调},
\end{cases}
\qquad(A2)
\end{align*}
其中当 $a>y_1$ 时区间 $[y_1-a,y_1+b]$ 按 $[0,y_1+b]$ 理解；
\begin{align*}
\begin{cases}
\phi\in C^\infty([0,\infty));\quad \lim_{y\to\infty}\phi(y)/y=0;\\
\text{存在 }\alpha\geqslant-1\text{ 与 }\beta\in\mathbb{R}\text{ 使对一切 }y\in[0,\infty),\\
\phi''(y)[\alpha\phi(y)+\beta]\leqslant-[\phi'(y)]^2<0.
\end{cases}
\qquad(A2')
\end{align*}

\begin{remark}\phantomsection\label{remark:ch07-7.1.1}
由 (A2')，映射 $y\mapsto\phi''(y)[\alpha\phi(y)+\beta]$ 在 $[0,\infty)$ 上严格负且连续，故只有两种可能：(a) $\alpha\phi+\beta>0$ 且 $\phi''<0$ 在 $[0,\infty)$ 上，或 (b) $\alpha\phi+\beta<0$ 且 $\phi''>0$ 在 $[0,\infty)$ 上。

在情形 (a)，$\phi''<0$ 使 $\phi'$ 在 $[0,\infty)$ 上单调递减，故 $\phi'(y)\leqslant\phi'(0)$ 对一切 $y\geqslant0$ 成立；又由 $\lim_{y\to\infty}\phi(y)/y=0$ 得 $\lim_{y\to\infty}\phi'(y)=0$（若 $\phi'$ 有非零极限 $L$，则 $\phi(y)/y\to L$），故 $\phi'(y)\geqslant0$。结合 (A2') 中 $-[\phi'(y)]^2<0$ 得 $0<\phi'\leqslant\phi'(0)$ 在 $[0,\infty)$ 上。类似地，在情形 (b)，$\phi''>0$ 使 $\phi'$ 单调递增，故 $\phi'(y)\geqslant\phi'(0)$；由 $\lim_{y\to\infty}\phi'(y)=0$ 得 $\phi'(y)\leqslant0$，结合 $[\phi'(y)]^2>0$ 得 $\phi'(0)\leqslant\phi'<0$ 在 $[0,\infty)$ 上。

特别地，两种情形下 (A2') 均蕴含 (A2) 与 (A1)：$\phi\in C^\infty([0,\infty))$ 给出 $\phi\in C^1([0,\infty))$，而 $\phi'$ 在整个 $[0,\infty)$ 上严格单调，故 (A2) 中对每个 $y_1>0$ 取 $a=b=1$ 即可；$\lim_{y\to\infty}\phi(y)/y=0$ 正是 (A1) 的第二条，而 $\phi\in C^\infty([0,\infty))$ 蕴含 $\phi\in\mathrm{Lip}([0,\infty))$。此外，若 $\phi(0)=0$，则 $\alpha\phi+\beta>0$ 在 $[0,\infty)$ 上当且仅当 $\beta>0$：在 $y=0$ 处 $\alpha\phi(0)+\beta=\beta$，故必要性成立；反之若 $\beta>0$，则 $\alpha\phi+\beta$ 在 $[0,\infty)$ 上不变号（否则由介值定理它在某点为零，与 $\phi''[\alpha\phi+\beta]\leqslant-[\phi']^2<0$ 矛盾），故其恒正。
\end{remark}

% SUPPLEMENT: parent=definition:ch07-流类; kind=motivation; location=7.1.1 小节 Remark 7.1.1 之后的译本定义及注记; basis=(A2') 原文的整理：把该结构假设提升为具名流类，以便引言中的引用有真实对象；两条情形的划分与具体成员 $\phi(y)=1-(1+y)^p$（$p<0$）的验证
\begin{definition}[正则流类]\label{definition:ch07-流类}
设 $\phi\in C^\infty([0,\infty))$ 且 $\lim_{y\to\infty}\phi(y)/y=0$。若存在常数 $\alpha\geqslant-1$ 与 $\beta\in\mathbb{R}$，使
\begin{align*}
\phi''(y)[\alpha\phi(y)+\beta]\leqslant-[\phi'(y)]^2<0\qquad\text{对一切 }y\in[0,\infty),
\end{align*}
则称 $\phi$ 属于\emph{正则流类}，并称满足上式的 $(\alpha,\beta)$ 为该流类的一个结构参数对。
\end{definition}

\begin{remark}\phantomsection\label{remark:ch07-流类动机}
(i) 动机。命题性结果（\cref{theorem:ch07-存在唯一}）只在满足 (A2') 的流上成立，故把 (A2') 所刻画的一类流取一个名称，可使"所构造熵解在何种流上满足相容条件并因而唯一"这一陈述有明确的对象。该类中的每个流都使 $\phi'$ 在 $[0,\infty)$ 上不变号且严格单调，这正是后文比较原理与正则性估计所依赖的结构。

(ii) 结构参数的作用。$\alpha\geqslant-1$ 使 $\alpha\phi+\beta$ 可以取到 $\beta$ 以外的一切符号，从而统一覆盖 $\phi''<0$（情形 (a)）与 $\phi''>0$（情形 (b)）两类流；$\beta$ 是 $\alpha\phi+\beta$ 在 $y=0$ 处的值，当 $\phi(0)=0$ 时其符号决定情形 (a) 与 (b)（\cref{remark:ch07-7.1.1}）。$[\phi'(y)]^2>0$ 是非退化条件，它排除 $\phi'$ 在某个区间上恒为零的退化情形。

(iii) 边界与例。该类非空：对 $p<0$ 与 $\phi(y)=1-(1+y)^p$，取 $\alpha=p/(1-p)$、$\beta=-p/(1-p)$，则 $\alpha\phi+\beta=-(p/(1-p))(1+y)^p>0$、$\phi''=-p(p-1)(1+y)^{p-2}<0$，直接计算给出 $\phi''[\alpha\phi+\beta]=-[\phi']^2$，即 (A2') 以等号成立（另见\cref{remark:ch07-7.2.2}）。反之，线性流 $\phi(y)=Cy$（$C\neq0$）不满足 (A2')，因 $\phi''\equiv0$ 使该不等式化为 $0<0$；此时 (P) 是线性输运方程，不需要该结构假设。
\end{remark}

流 $\phi$ 为超线性（形如 $\phi(y)=y^p$，$p>1$）时的问题 (P) 已在文献~\cite{ref67}中研究，其中证明了非负熵解的存在唯一性（另见文献~\cite{ref33}）。按该文定义，对正的时刻解取 $L^1(\mathbb{R})$ 中的值，尽管 $u_0$ 是有限 Radon 测度。关于 $\phi$ 在无穷远处至多线性增长时 (P) 的有趣结果，见开创性论文文献~\cite{ref39}。

当 $\phi(y)=Cy$（$C\in\mathbb{R}$）时，问题 (P) 即线性输运方程的 Cauchy 问题
\begin{align*}
\begin{cases}
u_t+Cu_x=0 & \text{在 }S,\\
u=u_0 & \text{在 }\mathbb{R}\times\{0\},
\end{cases}
\end{align*}
其解平凡地为 $u_0$ 沿直线 $x=Ct+x_0$（$x_0\in\mathbb{R}$）的平移。特别地，解的奇异部分 $u_s(\cdot,t)$ 对 $t>0$ 非零当且仅当它在 $t=0$ 时非零。

自然要问：当 $\phi$ 为次线性时会发生什么，若干数学模型正是如此提示（见文献~\cite{ref50,ref88,ref89}）。此情形下求 (P) 解的自然空间为 $L^\infty_{w^*}(0,T;\mathscr{R}_f^+(\mathbb{R}))$，如下列定义所设。

\subsection{解的定义}\label{subsection:ch07-解的定义}

本节始终假设 (A0)--(A1) 成立。沿用\cref{chapter:ch06}的记号：对 $u\in L^\infty_{w^*}(0,T;\mathscr{R}_f(\mathbb{R}))$ 与几乎所有 $t\in(0,T)$，把 $\mathbb{R}$ 上的有限 Radon 测度 $u(\cdot,t)$ 分解为
\begin{align*}
u(\cdot,t)=u_r(\cdot,t)\,\mathrm{d}x+u_s(\cdot,t),
\end{align*}
其中 $u_r(\cdot,t)\in L^1(\mathbb{R})$ 为关于 Lebesgue 测度的绝对连续部分的密度，$u_s(\cdot,t)\in\mathscr{R}_f(\mathbb{R})$ 为奇异部分；对 $u_0$ 使用同样的记号 $u_{0r}$、$u_{0s}$。$\langle\cdot,\cdot\rangle$ 表示 $C_0(\mathbb{R})$ 与 $\mathscr{R}_f(\mathbb{R})$ 之间的对偶配对。$\mathcal{Y}(S;\mathbb{R})$ 表示 $S\times\mathbb{R}$ 上这样的 Young 测度全体：其到 $S$ 的投影等于 Lebesgue 测度 $\lambda_2|_S$（见\cref{definition:ch02-Young测度}）。

\begin{definition}\label{definition:ch07-young测度解}
设 (A0)--(A1) 成立。若二元组 $(u,\nu)$ 满足下列条件 (i)--(iii)，则称它为问题 (P) 的 \emph{Young 测度解}：
\begin{enumerate}[(i)]
\item $u\in L^\infty_{w^*}(0,T;\mathscr{R}_f^+(\mathbb{R}))$，$\nu\in\mathcal{Y}(S;\mathbb{R})$；
\item 对几乎所有 $(x,t)\in S$，$\operatorname{supp}\nu_{(x,t)}\subseteq[0,\infty)$，且
\begin{align}
u_r(x,t)=\int_{[0,\infty)}y\,\mathrm{d}\nu_{(x,t)}(y), \label{align:ch07-7.2}
\end{align}
其中 $\{\nu_{(x,t)}\}_{(x,t)\in S}$ 为 $\nu$ 的离散分解（见\cref{definition:ch05-离散分解定义}）；
\item 对每个满足 $\zeta(\cdot,T)=0$ 于 $\mathbb{R}$ 的 $\zeta\in C^1([0,T];C_c^1(\mathbb{R}))$，有
\begin{align}
&\iint_S[u_r\partial_t\zeta+\phi^*\partial_x\zeta]\,\mathrm{d}x\,\mathrm{d}t
+\int_0^T\langle u_s(\cdot,t),\partial_t\zeta(\cdot,t)\rangle\,\mathrm{d}t\notag\\
&=-\langle u_0,\zeta(\cdot,0)\rangle, \label{align:ch07-7.3}
\end{align}
其中
\begin{align}
\phi^*(x,t)\triangleq\int_{[0,\infty)}\phi(y)\,\mathrm{d}\nu_{(x,t)}(y)
\quad\text{对几乎所有 }(x,t)\in S. \label{align:ch07-7.4}
\end{align}
\end{enumerate}
\end{definition}

% SUPPLEMENT: parent=definition:ch07-young测度解; kind=motivation; location=定义 7.1.1 后的注记; basis=定义 7.1.1 及其各项假设在 Young 测度框架中的作用：$u$ 的弱*可测性、$\operatorname{supp}\nu\subseteq[0,\infty)$ 与正性、重心公式 (7.2)、弱等式 (7.3) 中 $\phi^*$ 与 $u_s$ 的分工
\begin{remark}\phantomsection\label{remark:ch07-young测度解动机}
(i) 动机。当 $\phi$ 次线性时，(P) 的解一般不再是函数，而是 Radon 测度；逼近序列 $\{u_n\}$ 在 $L^1$ 中未必收敛，其弱极限只记录平均值，不记录 $u_n$ 的振荡。Young 测度 $\nu$ 逐点记录 $u_n(x,t)$ 的极限分布，从而在取极限时把非线性量 $\phi(u_n)$ 的极限写成 $\phi^*=\int\phi\,\mathrm{d}\nu_{(x,t)}$，这正是 (7.3) 中 $\phi^*$ 的来源。

(ii) 各条件的作用。$u\in L^\infty_{w^*}(0,T;\mathscr{R}_f^+(\mathbb{R}))$ 保证对几乎所有 $t$ 测度 $u(\cdot,t)$ 有定义、$\|u(\cdot,t)\|_{\mathscr{R}_f(\mathbb{R})}$ 本性有界，并使绝对连续部分 $u_r$ 与奇异部分 $u_s$ 可逐时刻分解；正性使 $u_r\geqslant0$。$\nu\in\mathcal{Y}(S;\mathbb{R})$ 的投影条件使纤维 $\nu_{(x,t)}$ 为概率测度，故 (7.2) 的重心有意义且 $u_r$ 由 $\nu$ 唯一确定。$\operatorname{supp}\nu_{(x,t)}\subseteq[0,\infty)$ 编码解的非负性：它使 $\phi$ 只需在 $[0,\infty)$ 上有定义，也使 (7.2)、(7.4) 的积分限为 $[0,\infty)$。等\eqref{align:ch07-7.3} 是 (P) 的弱形式：其中 $u_r\partial_t\zeta$ 与 $\phi^*\partial_x\zeta$ 对应方程的两个通量项，而 $\langle u_s,\partial_t\zeta\rangle$ 是奇异部分对时间试验函数的贡献，它不能并入 $u_r$ 项，因为 $u_s$ 关于 Lebesgue 测度奇异。

(iii) 边界。条件 (ii) 只确定 $\nu$ 的重心，不唯一确定 $\nu$：若 $\nu_{(x,t)}$ 与 $\tilde\nu_{(x,t)}$ 有相同重心与相同 $\phi$ 积分，则二者给出同一 $(u_r,\phi^*)$，故 (7.3) 无法区分它们。因此 Young 测度解一般远不唯一，需要\cref{definition:ch07-young熵解}的熵条件从中挑选物理解。
\end{remark}

置 $H_+(y)\triangleq\chi_{(0,\infty)}(y)$，$H_-(y)\triangleq-\chi_{(-\infty,0)}(y)$，$\operatorname{sgn}(y)\triangleq H_+(y)+H_-(y)$（$y\in\mathbb{R}$）；并记 $[a]_+\triangleq\max\{a,0\}$、$[a]_-\triangleq\max\{-a,0\}$（$a\in\mathbb{R}$）。

\begin{definition}\label{definition:ch07-young熵解}
设 (A0)--(A1) 成立，$(u,\nu)$ 为 (P) 的 Young 测度解。若对每个满足 $\zeta\geqslant0$ 与 $\zeta(\cdot,T)=0$ 于 $\mathbb{R}$ 的 $\zeta\in C^1([0,T];C_c^1(\mathbb{R}))$ 以及每个 $l\in[0,\infty)$，有
\begin{align}
&\iint_S[E_l^*\partial_t\zeta+F_l^*\partial_x\zeta]\,\mathrm{d}x\,\mathrm{d}t
+\int_0^T\langle u_s(\cdot,t),\partial_t\zeta(\cdot,t)\rangle\,\mathrm{d}t\notag\\
&\geqslant-\int_{\mathbb{R}}E_l(u_{0r}(x))\zeta(x,0)\,\mathrm{d}x-\langle u_{0s},\zeta(\cdot,0)\rangle, \label{align:ch07-7.5}
\end{align}
其中 $E_l(y)\triangleq|y-l|$、$F_l(y)\triangleq\operatorname{sgn}(y-l)[\phi(y)-\phi(l)]$（$y\in[0,\infty)$），且
\begin{align*}
E_l^*\triangleq\int_{[0,\infty)}E_l(y)\,\mathrm{d}\nu_{(x,t)}(y),\qquad
F_l^*\triangleq\int_{[0,\infty)}F_l(y)\,\mathrm{d}\nu_{(x,t)}(y),
\end{align*}
则称 $(u,\nu)$ 为 (P) 的 \emph{Young 测度熵解}。

若把上述不等式中的 $E_l$ 换为 $E_l(y)\triangleq[y-l]_+$、$F_l$ 换为 $F_l(y)\triangleq H_+(y-l)[\phi(y)-\phi(l)]$（$y\in[0,\infty)$），并要求所得不等式对上述所有 $\zeta$ 与 $l$ 成立，则称 $(u,\nu)$ 为 (P) 的 \emph{Young 测度熵下解}；若换为 $E_l(y)\triangleq[y-l]_-$ 与 $F_l(y)\triangleq H_-(y-l)[\phi(y)-\phi(l)]$，则称 $(u,\nu)$ 为 (P) 的 \emph{Young 测度熵上解}。
\end{definition}

% SUPPLEMENT: parent=definition:ch07-young熵解; kind=motivation; location=定义 7.1.2 后的注记; basis=定义 7.1.2 的 Kruzkov 型熵条件、$l\in[0,\infty)$ 与 $\zeta\geqslant0$ 的作用，以及由 $E_l=[\,\cdot\,]_++[\,\cdot\,]_-$ 得到下解/上解的方式
\begin{remark}\phantomsection\label{remark:ch07-young熵解动机}
(i) 动机。弱等\eqref{align:ch07-7.3} 只表达守恒律，不排除非物理解；对守恒律，Kružkov 熵条件以全部凸熵 $E_l(y)=|y-l|$ 与相应熵通量 $F_l$ 挑选唯一物理解。(7.5) 是这一条件在 Young 测度框架下的形式：它同时约束 $u_r$（通过 $E_l(u_r)$、$F_l(u_r)$ 的重心）与 $\nu$ 所记录的振荡。

(ii) 各条件的作用。$l\in[0,\infty)$ 与 $\operatorname{supp}\nu\subseteq[0,\infty)$ 一致，使 $\phi(l)$ 有定义；若取 $l<0$，则 $E_l$、$F_l$ 在 $[0,\infty)$ 上不再是关于 $l$ 的适当平移，故条件只对非负 $l$ 陈述。$\zeta\geqslant0$ 使 (7.5) 成为单侧不等式（比较原理所需的方向）；$\zeta(\cdot,T)=0$ 使终时刻的项消失。$E_l$、$F_l$ 在 $\nu$ 下的积分 $E_l^*$、$F_l^*$ 是 $\nu$ 的纤维上定义的量，它们把逐点熵条件"提升"到 Young 测度层面。右端的 $E_l(u_{0r})$ 与 $\langle u_{0s},\zeta(\cdot,0)\rangle$ 分别来自初值的绝对连续部分与奇异部分。

(iii) 边界。熵解一般不唯一确定 $\nu$：若 $\phi'$ 在 $\nu_{(x,t)}$ 的支撑集上为常数，则不同的 $\nu$ 可以满足同一族熵不等式（见\hyperref[remark:ch07-7.5.2]{关于 $\phi'$ 零点的注记}与\hyperref[remark:ch07-7.5.3]{关于 Young 测度支撑的注记}）；\cref{theorem:ch07-7.3.2}表明在 (A2) 下 $\nu$ 退化为 Dirac 测度。把 (7.5) 与"用 $[\,\cdot\,]_+$、$[\,\cdot\,]_-$ 分别替换 $E_l$、$F_l$"相结合，得到熵下解与熵上解；由于 $E_l=[\,\cdot\,]_++[\,\cdot\,]_-$、$\operatorname{sgn}=H_++H_-$，熵解恰是二者之和所对应的条件。
\end{remark}

\begin{remark}\phantomsection\label{remark:ch07-7.1.2}
(i) 由\eqref{align:ch07-7.1}、\eqref{align:ch07-7.2}与\eqref{align:ch07-7.4}得
\begin{align}
|\phi^*(x,t)|\leqslant M\int_{[0,\infty)}y\,\mathrm{d}\nu_{(x,t)}(y)=Mu_r(x,t)
\quad\text{对几乎所有 }(x,t)\in S. \label{align:ch07-7.6}
\end{align}
因 $u_r\in L^\infty(0,T;L^1(\mathbb{R}))$，由\eqref{align:ch07-7.6}知 $\phi^*\in L^\infty(0,T;L^1(\mathbb{R}))$。

(ii) 估计\eqref{align:ch07-7.6}只用到 $\phi$ 的线性增长界与 $\nu$ 的支撑条件，不要求 $\phi$ 有界或单调；它是后文对 $\phi^*$ 取极限与估计时间正则性的基础。
\end{remark}

对任意非空开集 $\Omega\subseteq\mathbb{R}$ 与 $\tau\in(0,T)$，置 $Q_\tau\triangleq\Omega\times(0,\tau)$。

\begin{definition}\label{definition:ch07-熵解}
设 (A0)--(A1) 成立，$\Omega\subseteq\mathbb{R}$ 为非空开集，$\tau\in(0,T)$。若 $u\in L^\infty_{w^*}(0,T;\mathscr{R}_f^+(\Omega))$，且对每个满足 $\zeta(\cdot,\tau)=0$ 于 $\Omega$ 的 $\zeta\in C^1([0,\tau];C_c^1(\Omega))$，有
\begin{align}
&\iint_{Q_\tau}[u_r\partial_t\zeta+\phi(u_r)\partial_x\zeta]\,\mathrm{d}x\,\mathrm{d}t
+\int_0^\tau\langle u_s(\cdot,t),\partial_t\zeta(\cdot,t)\rangle\,\mathrm{d}t\notag\\
&=-\langle u_0,\zeta(\cdot,0)\rangle, \label{align:ch07-7.7}
\end{align}
则称 $u$ 为问题 (P) 在 $Q_\tau$ 中的\emph{解}。
\end{definition}

(P) 在 $S$ 中的解简称为 (P) 的解。

% SUPPLEMENT: parent=definition:ch07-熵解; kind=motivation; location=定义 7.1.3 后的注记; basis=定义 7.1.3 的弱形式、局部区域 $Q_\tau$ 的引入动机，以及 (7.7) 与 (7.3) 中 $\phi(u_r)$、$\phi^*$ 的差别
\begin{remark}\phantomsection\label{remark:ch07-熵解动机}
(i) 动机。本定义把"测度 $u$ 解 (P)"的直观说法写成只涉及试验函数的等式：$u_r$ 与 $u_s$ 分别测试 $\partial_t\zeta$，$\phi(u_r)$ 测试 $\partial_x\zeta$。与\cref{definition:ch07-young测度解}相比，这里不再有 Young 测度，$\phi(u_r)$ 直接由 $u$ 的绝对连续部分算出；这正是把 Young 测度解"投影"回单个测度后的弱形式。

(ii) 各条件的作用。$\Omega$ 与 $\tau$ 使等式可在局部区域 $Q_\tau$ 上陈述，从而允许只要求解在局部满足方程；$\zeta(\cdot,\tau)=0$ 去掉终时刻的边界项，使 (7.7) 对每个 $\tau$ 都成立而不引入 $\tau$ 处的未知值。$u\in L^\infty_{w^*}(0,T;\mathscr{R}_f^+(\Omega))$ 保证 $u_r$、$u_s$ 逐时刻有定义且 $\phi(u_r)$ 可积（由 (7.1) 的 $|\phi(u_r)|\leqslant Mu_r$）。

(iii) 边界。$Q_\tau$ 中的解不必是 $S$ 中的解：后者要求 $\Omega=\mathbb{R}$ 且等式对一切 $\tau\in(0,T)$ 与一切 $\zeta$ 成立。反过来，$S$ 中的解限制到 $Q_\tau$ 上满足 (7.7)（取试验函数支集于 $\Omega$ 即可），故局部化不改变解类。\eqref{align:ch07-7.7}中的 $\phi(u_r)$ 与\eqref{align:ch07-7.3}中的 $\phi^*$ 一般不同：前者只依赖 $u$，后者依赖 Young 测度 $\nu$。
\end{remark}

\begin{definition}\label{definition:ch07-Qtau熵解}
设 (A0)--(A1) 成立，$\Omega\subseteq\mathbb{R}$ 为非空开集，$\tau\in(0,T)$，$u$ 为 (P) 在 $Q_\tau$ 中的解。若对所有满足 $\zeta\geqslant0$ 与 $\zeta(\cdot,\tau)=0$ 于 $\Omega$ 的 $\zeta\in C^1([0,\tau];C_c^1(\Omega))$ 以及所有 $l\in[0,\infty)$，有
\begin{align}
&\iint_{Q_\tau}\{E_l(u_r)\partial_t\zeta+F_l(u_r)\partial_x\zeta\}\,\mathrm{d}x\,\mathrm{d}t
+\int_0^\tau\langle u_s(\cdot,t),\partial_t\zeta(\cdot,t)\rangle\,\mathrm{d}t\notag\\
&\geqslant-\int_\Omega E_l(u_{0r}(x))\zeta(x,0)\,\mathrm{d}x
-\langle u_{0s},\zeta(\cdot,0)\rangle, \label{align:ch07-7.8}
\end{align}
则称 $u$ 为 (P) 在 $Q_\tau$ 中的\emph{熵解}，并称\eqref{align:ch07-7.8}为熵不等式。
\end{definition}

% SUPPLEMENT: parent=definition:ch07-Qtau熵解; kind=motivation; location=定义 7.1.4 后的注记; basis=定义 7.1.4 的熵不等式：只对非负 $l$ 与 $\zeta\geqslant0$ 要求、$E_l$ 与 $F_l$ 的凸性含义，以及它与定义 7.1.2 的对应关系
\begin{remark}\phantomsection\label{remark:ch07-Qtau熵解动机}
(i) 动机。本定义是 Young 测度熵解的定义在单测度情形下的对应物：把 $E_l^*$、$F_l^*$ 换成 $E_l(u_r)$、$F_l(u_r)$，即把 Young 测度上的积分换成由 $u$ 的绝对连续部分直接算出的值。它是唯一性证明（\cref{section:ch07-唯一性}）中实际使用的条件。

(ii) 各条件的作用。$l\in[0,\infty)$ 使 $\phi(l)$ 有定义并与非负性相容；$\zeta\geqslant0$ 给出不等式的方向；$\zeta(\cdot,\tau)=0$ 去掉终时刻项；$E_l(y)=|y-l|$ 是 Kružkov 熵，$F_l$ 是相应熵通量，二者的凸性使 (7.8) 对 $l$ 的积分与极限运算封闭。奇异部分以 $\langle u_s,\partial_t\zeta\rangle$ 的形式出现，与\cref{definition:ch07-young测度解}中的处理一致。

(iii) 边界。对 $l=0$，(7.8) 化为关于 $u_r$ 的正性型不等式；对一般的 $l$，(7.8) 蕴含 $L^1$ 压缩性。若不要求 $\zeta\geqslant0$，则不等式不再有比较意义：取 $-\zeta$ 会反转方向。
\end{remark}

\begin{definition}\label{definition:ch07-上下解}
设 (A0)--(A1) 成立，$\Omega\subseteq\mathbb{R}$ 为非空开集，$\tau\in(0,T)$。若 $u\in L^\infty_{w^*}(0,T;\mathscr{R}_f^+(\Omega))$ 且对所有满足 $\zeta\geqslant0$ 与 $\zeta(\cdot,\tau)=0$ 于 $\Omega$ 的 $\zeta\in C^1([0,\tau];C_c^1(\Omega))$ 以及所有 $l\in[0,\infty)$，有
\begin{align}
&\iint_{Q_\tau}\{[u_r-l]_+\partial_t\zeta+H_+(u_r-l)[\phi(u_r)-\phi(l)]\partial_x\zeta\}\,\mathrm{d}x\,\mathrm{d}t\notag\\
&\quad+\int_0^\tau\langle u_s(\cdot,t),\partial_t\zeta(\cdot,t)\rangle_\Omega\,\mathrm{d}t
\geqslant-\int_\Omega[u_{0r}-l]_+\zeta(x,0)\,\mathrm{d}x
-\langle u_{0s},\zeta(\cdot,0)\rangle_\Omega, \label{align:ch07-7.9}
\end{align}
则称 $u$ 为 (P) 在 $Q_\tau$ 中的\emph{熵下解}；若对所有上述 $\zeta$ 与 $l$，有
\begin{align}
&\iint_{Q_\tau}\{[u_r-l]_-\partial_t\zeta+H_-(u_r-l)[\phi(u_r)-\phi(l)]\partial_x\zeta\}\,\mathrm{d}x\,\mathrm{d}t\notag\\
&\geqslant-\int_\Omega[u_{0r}-l]_-\zeta(x,0)\,\mathrm{d}x, \label{align:ch07-7.10}
\end{align}
则称 $u$ 为 (P) 在 $Q_\tau$ 中的\emph{熵上解}。
\end{definition}

% SUPPLEMENT: parent=definition:ch07-上下解; kind=motivation; location=定义 7.1.5 后的注记; basis=定义 7.1.5 的两条不等式及"每个熵解既是熵下解又是熵上解"的代数验证：利用 (7.7) 消去线性部分，说明 (7.9) 与 (7.10) 各自等价于 (7.8)
\begin{remark}\phantomsection\label{remark:ch07-上下解动机}
(i) 动机。熵下解与熵上解把 (7.8) 按 $E_l=[\,\cdot\,]_++[\,\cdot\,]_-$ 与 $\operatorname{sgn}=H_++H_-$ 拆成两半，分别对应 $u_r$ 位于 $l$ 上方与下方的部分。它们是比较原理与唯一性证明（\cref{lemma:ch07-7.2.3}、\cref{theorem:ch07-唯一性}）的直接工具：只需对熵上解证明单侧估计，再对 $u-v$ 应用。

(ii) 各条件的作用。$\zeta\geqslant0$ 与 $l\in[0,\infty)$ 的作用与熵解定义中的相同。$H_+$ 与 $H_-$ 分别把熵通量限制在 $u_r>l$ 与 $u_r<l$ 的区域，这正是 $[\,\cdot\,]_\pm$ 的导数；$\langle u_s,\partial_t\zeta\rangle_\Omega$ 出现在 (7.9) 中，而 (7.10) 中不出现该项。

(iii) 边界与两条不等式的关系。设 $u$ 为 (P) 在 $Q_\tau$ 中的熵解。把 $[u_r-l]_\pm=\frac12(|u_r-l|\pm(u_r-l))$ 与 $H_\pm=\frac12(\operatorname{sgn}\pm1)$ 代入，并利用\eqref{align:ch07-7.7}消去线性部分（注意 $\int_0^\tau\int_\Omega\partial_t\zeta\,\mathrm{d}x\,\mathrm{d}t=-\int_\Omega\zeta(x,0)\,\mathrm{d}x$、$\int_0^\tau\int_\Omega\partial_x\zeta\,\mathrm{d}x\,\mathrm{d}t=0$），可将 (7.9) 与 (7.10) 的左端分别化为 $\frac12$(7.8)左端 $+\frac12(-\langle u_0,\zeta(\cdot,0)\rangle+l\int_\Omega\zeta(x,0)\,\mathrm{d}x)$ 与 $\frac12$(7.8)左端 $+\frac12(\langle u_0,\zeta(\cdot,0)\rangle-l\int_\Omega\zeta(x,0)\,\mathrm{d}x)$，而右端相应化为 $\frac12$(7.8)右端加上同一个量（对 (7.9) 还多出 $-\frac12\langle u_{0s},\zeta(\cdot,0)\rangle_\Omega$）。因此 (7.10) 与 (7.8) 等价，(7.9) 由 (7.8) 与 $\langle u_{0s},\zeta(\cdot,0)\rangle_\Omega\geqslant0$ 推出。这解释了为何 (7.10) 中不出现奇异项：它与线性部分的贡献恰好相消。特别地，每个熵解既是熵下解又是熵上解；反之，把 (7.9) 与 (7.10) 相加（并用 $u_{0s}\geqslant0$）即得 (7.8)，故两个类之交恰为熵解类。
\end{remark}

\begin{remark}\phantomsection\label{remark:ch07-7.1.3}
(i) 等式\eqref{align:ch07-7.7}亦可写为
\begin{align*}
\langle u,\partial_t\zeta\rangle+\iint_S\phi(u_r)\partial_x\zeta\,\mathrm{d}x\,\mathrm{d}t=-\langle u_0,\zeta(\cdot,0)\rangle,
\end{align*}
表明 $u_t=-[\phi(u_r)]_x$ 在 $\mathcal{D}^*(S)$ 中。

(ii) 问题 (P) 的解也是 Young 测度解。此外由\eqref{align:ch07-7.1}知 $\phi(u_r)\in L^\infty(0,T;L^1(\mathbb{R}))$。熵解、熵下解与熵上解有类似注记。
\end{remark}

下列结果与\cref{proposition:ch6-1.1}类似。其证明亦类似，故从略（见文献~\cite{ref17} 命题 3.5）；下面给出完整论证。

\begin{proposition}\label{proposition:ch07-7.1.1}
设 (A0)--(A1) 成立，$(u,\nu)$ 为 (P) 的 Young 测度解，$\phi^*$ 由\eqref{align:ch07-7.4}给出。则：
\begin{enumerate}[(i)]
\item 存在零测集 $F^*\subseteq(0,T)$，使对一切 $t\in(0,T)\setminus F^*$ 与一切 $\rho\in C_c^1(\mathbb{R})$，有
\begin{align}
\langle u(\cdot,t),\rho\rangle-\langle u_0,\rho\rangle
=\int_0^t\int_{\mathbb{R}}\phi^*(x,s)\rho'(x)\,\mathrm{d}x\,\mathrm{d}s; \label{align:ch07-7.11}
\end{align}
其中 $\langle u(\cdot,t),\rho\rangle\triangleq\int_{\mathbb{R}}\rho u_r(\cdot,t)\,\mathrm{d}x+\langle u_s(\cdot,t),\rho\rangle$；
\item 对一切 $\rho\in C_c(\mathbb{R})$，有
\begin{align}
\operatorname*{ess\,lim}_{t\to0^+}\langle u(\cdot,t),\rho\rangle=\langle u_0,\rho\rangle, \label{align:ch07-7.12}\\
\operatorname*{ess\,lim}_{t\to t_0}\langle u(\cdot,t),\rho\rangle=\langle u(\cdot,t_0),\rho\rangle
\quad\text{对几乎所有 }t_0\in(0,T). \label{align:ch07-7.13}
\end{align}
\end{enumerate}
\end{proposition}

% SUPPLEMENT: parent=proposition:ch07-7.1.1; kind=proof; location=命题 7.1.1 之后的完整证明（原书仅注明"证明类似，从略"）; basis=由 (7.3) 取乘积试验函数得到 $\langle u(\cdot,t),\rho\rangle$ 的分布导数，用绝对连续代表元与可数稠密试验函数族证明 (7.11)--(7.13)；并逐项核对 (7.1) 给出的 $\phi^*\in L^\infty(0,T;L^1(\mathbb{R}))$
\begin{proof}
记 $M_0\triangleq\operatorname*{ess\,sup}_{t\in(0,T)}\|u(\cdot,t)\|_{\mathscr{R}_f(\mathbb{R})}<\infty$。设 $\rho\in C_c^1(\mathbb{R})$，对 $t\in(0,T)$ 置
\begin{align*}
g_\rho(t)&\triangleq\langle u(\cdot,t),\rho\rangle
=\int_{\mathbb{R}}\rho(x)u_r(x,t)\,\mathrm{d}x+\langle u_s(\cdot,t),\rho\rangle,\\
h_\rho(t)&\triangleq\int_{\mathbb{R}}\phi^*(x,t)\rho'(x)\,\mathrm{d}x.
\end{align*}
由 $u\in L^\infty_{w^*}(0,T;\mathscr{R}_f^+(\mathbb{R}))$ 得 $|g_\rho(t)|\leqslant M_0\|\rho\|_\infty$ 对几乎所有 $t$ 成立；由\eqref{align:ch07-7.6}得
\begin{align*}
|h_\rho(t)|\leqslant\|\phi^*(\cdot,t)\|_{L^1(\mathbb{R})}\|\rho'\|_\infty
\leqslant M\|u_r(\cdot,t)\|_{L^1(\mathbb{R})}\|\rho'\|_\infty
\leqslant MM_0\|\rho'\|_\infty,
\end{align*}
故 $g_\rho,h_\rho\in L^\infty(0,T)$。

设 $\beta\in C^1([0,T])$ 满足 $\beta(T)=0$。则 $\zeta(x,t)\triangleq\rho(x)\beta(t)$ 属于 $C^1([0,T];C_c^1(\mathbb{R}))$ 且 $\zeta(\cdot,T)=0$，故可代入\eqref{align:ch07-7.3}；注意 $\partial_t\zeta=\rho\beta'$、$\partial_x\zeta=\rho'\beta$ 以及 $\langle u_s(\cdot,t),\rho\beta'(t)\rangle=\beta'(t)\langle u_s(\cdot,t),\rho\rangle$，得
\begin{align}
\int_0^Tg_\rho(t)\beta'(t)\,\mathrm{d}t+\int_0^Th_\rho(t)\beta(t)\,\mathrm{d}t
=-\langle u_0,\rho\rangle\beta(0). \label{align:ch07-7.1proofa}
\end{align}
置 $G_\rho(t)\triangleq\langle u_0,\rho\rangle+\int_0^th_\rho(s)\,\mathrm{d}s$（$t\in[0,T]$）。则 $G_\rho$ 绝对连续，$G_\rho(0)=\langle u_0,\rho\rangle$，且 $G_\rho'=h_\rho$ 几乎处处。由分部积分，
\begin{align*}
\int_0^TG_\rho(t)\beta'(t)\,\mathrm{d}t
=-\langle u_0,\rho\rangle\beta(0)-\int_0^Th_\rho(t)\beta(t)\,\mathrm{d}t .
\end{align*}
与\eqref{align:ch07-7.1proofa}相减，得
\begin{align*}
\int_0^T(g_\rho-G_\rho)(t)\beta'(t)\,\mathrm{d}t=0
\end{align*}
对一切满足 $\beta(T)=0$ 的 $\beta\in C^1([0,T])$ 成立。先取 $\beta\in C_c^\infty(0,T)$，得 $(g_\rho-G_\rho)'=0$ 于 $\mathcal{D}'(0,T)$，故存在常数 $c_\rho$ 使 $g_\rho-G_\rho=c_\rho$ 几乎处处；再取 $\beta(t)\triangleq(1-t/T)_+$（它属于 $C^1([0,T])$、$\beta(T)=0$），得 $c_\rho(\beta(T)-\beta(0))=-c_\rho=0$，即 $c_\rho=0$。于是
\begin{align}
g_\rho(t)=\langle u_0,\rho\rangle+\int_0^th_\rho(s)\,\mathrm{d}s
\quad\text{对几乎所有 }t\in(0,T), \label{align:ch07-7.1proofb}
\end{align}
这正是\eqref{align:ch07-7.11}对 $\rho\in C_c^1(\mathbb{R})$ 的情形。

现取可数集 $D\subseteq C_c^1(\mathbb{R})$，使对每个 $\rho\in C_c^1(\mathbb{R})$ 与每个 $\varepsilon>0$，存在 $\sigma\in D$ 满足 $\|\rho-\sigma\|_\infty+\|\rho'-\sigma'\|_\infty<\varepsilon$（例如取 $D$ 为具有有理系数、支集端点与取值均为有理数的 $C_c^\infty(\mathbb{R})$ 函数之全体）。对每个 $\sigma\in D$，由\eqref{align:ch07-7.1proofb}存在零测集 $N_\sigma\subseteq(0,T)$，使 $t\in(0,T)\setminus N_\sigma$ 时\eqref{align:ch07-7.11}对 $\sigma$ 成立。置 $F^*\triangleq\bigcup_{\sigma\in D}N_\sigma$，则 $F^*$ 为零测集。设 $t\in(0,T)\setminus F^*$ 与 $\rho\in C_c^1(\mathbb{R})$，取 $\{\sigma_k\}\subseteq D$ 使 $\|\rho-\sigma_k\|_\infty+\|\rho'-\sigma_k'\|_\infty\to0$。由
\begin{align*}
|g_\rho(t)-g_{\sigma_k}(t)|\leqslant M_0\|\rho-\sigma_k\|_\infty
\end{align*}
以及
\begin{align*}
&\left|\Big(\langle u_0,\rho\rangle+\int_0^t\int_{\mathbb{R}}\phi^*\rho'\Big)
-\Big(\langle u_0,\sigma_k\rangle+\int_0^t\int_{\mathbb{R}}\phi^*\sigma_k'\Big)\right|\\
&\leqslant\|u_0\|_{\mathscr{R}_f(\mathbb{R})}\|\rho-\sigma_k\|_\infty
+t\sup_{s\in(0,T)}\|\phi^*(\cdot,s)\|_{L^1(\mathbb{R})}\|\rho'-\sigma_k'\|_\infty,
\end{align*}
令 $k\to\infty$ 即知\eqref{align:ch07-7.11}对 $\rho$ 与 $t$ 成立。于是证明了 (i)。

(ii) 设 $\rho\in C_c^1(\mathbb{R})$。由\eqref{align:ch07-7.1proofb}，$g_\rho$ 几乎处处等于连续函数 $G_\rho$，故
\begin{align*}
\operatorname*{ess\,lim}_{t\to0^+}g_\rho(t)=G_\rho(0)=\langle u_0,\rho\rangle,
\end{align*}
且对几乎所有 $t_0\in(0,T)$（即那些使 $g_\rho(t_0)=G_\rho(t_0)$ 的 $t_0$）有 $\operatorname*{ess\,lim}_{t\to t_0}g_\rho(t)=G_\rho(t_0)=g_\rho(t_0)$。

设 $\rho\in C_c(\mathbb{R})$。取 $\{\rho_k\}\subseteq C_c^1(\mathbb{R})$ 与紧集 $K\subseteq\mathbb{R}$，使 $\operatorname{supp}\rho_k\subseteq K$ 且 $\|\rho_k-\rho\|_\infty\to0$。则对几乎所有 $t$ 有 $|g_\rho(t)-g_{\rho_k}(t)|\leqslant M_0\|\rho-\rho_k\|_\infty$，且 $|\langle u_0,\rho\rangle-\langle u_0,\rho_k\rangle|\leqslant\|u_0\|_{\mathscr{R}_f(\mathbb{R})}\|\rho-\rho_k\|_\infty$。把上述结论对每个 $\rho_k$ 写出，取相应的例外零测集之并（仍为零测集），再令 $k\to\infty$，即得\eqref{align:ch07-7.12}与\eqref{align:ch07-7.13}对 $\rho$ 成立。结论得证。
\end{proof}

% SUPPLEMENT: parent=proposition:ch07-7.1.1; kind=motivation; location=命题 7.1.1 证明后的注记; basis=命题 7.1.1 的两项结论：初值在弱*意义下被取到、$t\mapsto u(\cdot,t)$ 的代表元的弱*连续性与绝对连续性；逐项分析 (A1) 与 Young 测度解定义的作用及零测集的必要性
\begin{remark}\phantomsection\label{remark:ch07-7.1.1动机}
(i) 动机。Young 测度解定义中的弱等式\eqref{align:ch07-7.3}只说明 $u$ 在分布意义下满足方程，既未说明初值 $u_0$ 以何种意义被取到，也未说明映射 $t\mapsto u(\cdot,t)$ 有任何连续性。本命题补足这两点：\eqref{align:ch07-7.11}把 $\langle u(\cdot,t),\rho\rangle-\langle u_0,\rho\rangle$ 写成 $\int_0^t\int\phi^*\rho'$，故 $g_\rho$ 与一个绝对连续函数几乎处处相等；\eqref{align:ch07-7.12}说明初值在弱*意义下被取到；\eqref{align:ch07-7.13}说明 $g_\rho$ 几乎处处连续。这些事实在\hyperref[remark:ch07-7.1.4]{关于连续代表元的注记}与唯一性论证中反复使用。

(ii) 各假设的作用。(A1) 通过\eqref{align:ch07-7.6}给出 $\phi^*\in L^\infty(0,T;L^1(\mathbb{R}))$，这是 $h_\rho\in L^\infty(0,T)$、从而 $g_\rho$ 绝对连续的唯一依据；它还保证 $|\phi(u_r)|\leqslant Mu_r$，使 $\phi(u_r)\in L^\infty(0,T;L^1(\mathbb{R}))$。$u\in L^\infty_{w^*}(0,T;\mathscr{R}_f^+(\mathbb{R}))$ 给出与 $t$ 无关的范数界 $M_0$，它在把\eqref{align:ch07-7.11}从可数稠密集 $D$ 过渡到一切 $\rho\in C_c^1(\mathbb{R})$ 时控制误差。$\nu$ 只通过 $\phi^*$ 进入证明，故结论对一切 Young 测度解成立而不依赖 $\nu$ 的具体结构。$\rho\in C_c^1(\mathbb{R})$ 使 $\rho'$ 有界且支集紧，$\rho\in C_c(\mathbb{R})$ 时的结论由一致逼近得到。

(iii) 边界。零测集 $F^*$ 不能去掉：$t\mapsto\langle u(\cdot,t),\rho\rangle$ 只是 $L^\infty$ 函数，在零测集上可任意改动。\eqref{align:ch07-7.12}只是弱*收敛而非范数收敛：当 $u_0$ 含 Dirac 质量时，$\|u(\cdot,t)-u_0\|_{\mathscr{R}_f(\mathbb{R})}$ 一般不必趋于零（见\hyperref[remark:ch07-7.2.1]{关于等待时间的注记}）。若去掉 (A1) 中 $\lim_{y\to\infty}\phi(y)/y=0$ 而允许超线性增长，则\eqref{align:ch07-7.6}失效，$\phi^*$ 未必属于 $L^\infty(0,T;L^1(\mathbb{R}))$，\eqref{align:ch07-7.11}右端的积分可能不绝对收敛。
\end{remark}

\begin{remark}\phantomsection\label{remark:ch07-7.1.4}
如\hyperref[remark:ch06-6.1.2]{关于连续代表元的注记}，由\cref{proposition:ch07-7.1.1}易见映射 $t\mapsto u(\cdot,t)$ 有一个对任意 $t\in[0,T]$ 都有定义的代表元，使
\begin{align}
\lim_{t\to t_0}\langle u(\cdot,t),\rho\rangle=\langle u(\cdot,t_0),\rho\rangle
\quad\text{对所有 }t_0\in[0,T]\text{ 与 }\rho\in C_c(\mathbb{R}). \label{align:ch07-7.14}
\end{align}
今后凡对 $t$ 在 $[0,T]$ 的某子区间上陈述映射 $t\mapsto u(\cdot,t)$ 的性质时，均指此连续代表元。
\end{remark}

如\cref{chapter:ch06}所处理的抛物问题（见\cref{section:ch06-唯一性}），正则性信息对证明唯一性很重要。这方面的首个结果部分地使人想起不等式\eqref{align:ch6-6.33}，陈述如下；其证明见\cref{section:ch07-正则证明}。

\begin{proposition}\label{proposition:ch07-7.1.2}
设 (A0)--(A1) 成立，$u\in L^\infty_{w^*}(0,T;\mathscr{R}_f^+(\mathbb{R}))$。
\begin{enumerate}[(i)]
\item 设 $u$ 为 (P) 的 Young 测度熵解（见\cref{definition:ch07-young熵解}）。则
\begin{align}
u_s(\cdot,t_2)\leqslant u_s(\cdot,t_1)\quad\text{在 }\mathscr{R}_f^+(\mathbb{R})\text{ 中对几乎所有 }0\leqslant t_1\leqslant t_2\leqslant T. \label{align:ch07-7.15}
\end{align}
特别地，
\begin{align}
u_s(\cdot,t)\leqslant u_{0s}\quad\text{在 }\mathscr{R}_f^+(\mathbb{R})\text{ 中对几乎所有 }t\in(0,T). \label{align:ch07-7.16}
\end{align}
\item 设 $u$ 为问题 (P) 的解（见\cref{definition:ch07-熵解}）。则质量守恒成立：
\begin{align}
\|u(\cdot,t)\|_{\mathscr{R}_f(\mathbb{R})}=\|u_0\|_{\mathscr{R}_f(\mathbb{R})}
\quad\text{对几乎所有 }t\in(0,T). \label{align:ch07-7.17}
\end{align}
\end{enumerate}
\end{proposition}

% SUPPLEMENT: parent=proposition:ch07-7.1.2; kind=motivation; location=命题 7.1.2 后的注记; basis=命题 7.1.2 的动机（正则性、奇异部分的单调性、质量守恒）与两项结论各自假设的作用；证明位置在后续小节
\begin{remark}\phantomsection\label{remark:ch07-7.1.2动机}
(i) 动机。唯一性证明需要两类信息：其一，奇异部分 $u_s(\cdot,t)$ 关于 $t$ 单调不增（\eqref{align:ch07-7.15}），它使相容条件所涉及的原子质量只能减小，从而可以把解按原子位置分段比较；其二，正则化信息，即解在正时刻的部分性质更好，它使 (P) 的解在 $t>0$ 时可与 Young 测度熵解比较。质量守恒\eqref{align:ch07-7.17}则是守恒律的基本性质，也用于把弱*收敛提升为范数收敛。

(ii) 各假设的作用。(A0)--(A1) 保证 $\phi$ 至多线性增长，从而 $\phi(u_r)\in L^\infty(0,T;L^1(\mathbb{R}))$；这是证明中对熵通量取极限的基础。(i) 只假设 $u$ 是 Young 测度熵解，因为证明只用到熵不等式\eqref{align:ch07-7.5}在 $l\to\infty$ 时的形式；(ii) 则假设 $u$ 是通常意义下的解，因为质量守恒要利用\eqref{align:ch07-7.7}中的线性结构。

(iii) 边界。两式都只对几乎所有 $t$ 成立，零测集一般依赖 $u$ 而不能去掉。\eqref{align:ch07-7.15}是测度意义下的不等式（对每个非负试验函数配对后成立），不是密度的逐点不等式；它允许 $u_s$ 在零测集上取任意值。若去掉正性假设（例如考虑带号初值），\eqref{align:ch07-7.16}不再有意义。
\end{remark}

\end{document}
